% Applied Energy submission -- Elsevier elsarticle class
% Time Series 2026: weighted representative energy extreme events and their change under climate change
% Template: minimal copy of HPP_potential_2026/PUBLICATION/Applied_Energy_2026 (2026-11-20)
% Content follows 01_research_plan.md; write into the sections as the project advances.

\documentclass[preprint,12pt]{elsarticle}

% Packages
\usepackage[T1]{fontenc}
\usepackage{amssymb}
\usepackage{amsmath}
\usepackage{graphicx}
\usepackage{booktabs}
\usepackage{array}
\usepackage{multirow}
\usepackage{url}
\usepackage{xurl}
\usepackage{lineno}
\usepackage{siunitx}
\sisetup{detect-all,group-separator={\,}}

% Line numbers for review (comment out for the final version)
\linenumbers

% Auto-generated \newcommand values from the analysis code
% Single source of truth: do not edit the values by hand. Re-run the generating script to refresh.
% Auto-generated \newcommand values for numbers cited in the text.
% Do not edit by hand: write the generating script (reads the analysis CSVs) and re-run it.
% No values yet (2026-11-20).


% Marker for results that wait for the analysis (draft stage; remove before submission)
\usepackage{xcolor}
\newcommand{\Pending}[1]{\textcolor{red}{[Pending: #1]}}

\journal{Applied Energy}

\begin{document}

\begin{frontmatter}

% TITLE TO BE DEFINED: working title from the research topic (01_research_plan.md section 1)
\title{Weighted Representative Energy Extreme Events for Adequacy-Oriented Planning
and Their Representativeness under Climate Change}

% AUTHOR LIST AND ORDER TO BE DEFINED by the first author.
\author[erg]{Chao Tang\corref{cor1}}
\ead{chao.tang@univ-reunion.fr}
\cortext[cor1]{Corresponding author}

\author[dtu]{Matti Koivisto}

\affiliation[erg]{organization={ENERGY-Lab, University of Reunion},
            addressline={15 avenue Ren\'{e} Cassin, CS 92003},
            postcode={97744},
            city={Saint-Denis Cedex 9},
            country={Reunion, France}}

\affiliation[dtu]{organization={DTU Wind and Energy Systems, Technical University of Denmark},
            city={Roskilde},
            country={Denmark}}

\begin{abstract}
% TODO: write once the results are final (at most 250 words).
% Structure: context / gap / data and methods / main findings / implication.
\Pending{abstract}
\end{abstract}

\begin{highlights}
% 3--5 bullets, each at most 85 characters including spaces.
\item \Pending{highlight 1}
\item \Pending{highlight 2}
\item \Pending{highlight 3}
\end{highlights}

\begin{keyword}
% 1--7 keywords; to be finalised.
Energy extreme events \sep Resource adequacy \sep Time series aggregation \sep
Climate change \sep Residual load \sep Energy system planning
\end{keyword}

\end{frontmatter}

\section{Introduction}
\label{sec:introduction}
% Outline (01_research_plan.md sections 1-2); prose to be written.
% - Context: weather-driven systems need adequacy-oriented inputs; second stage of two-stage planning.
% - Gap: system-defined events are not comparable across years; resource-based indices omit demand
%   and the system configuration; no study tests a representative event set under climate change.
% - Weather extremes are not energy extremes: they depend on system size, capacity and interconnection.
% - Aim and research questions Q1-Q4.
\Pending{introduction}

\section{Data}
\label{sec:data}
% Outline (01_research_plan.md section 6; Time_Series_2026_DATA.md).
% - Historical: DTU Data collection doi:10.11583/DTU.c.7964141 (Balmorel regions, ERA5 1980/1982-2021):
%   electricity demand, wind, solar PV, hydro inflow.
% - Future: Mopo climate projections (wind, PV); PECD 4.2 (temperature, hydro); demand method.
% - Why this data set (01_research_plan.md D16, D17): the collection provides the only hourly,
%   multi-country electricity demand found that is consistent with the ERA5 weather years; wind, solar PV and
%   hydro inflow come from the same collection, on the same regions and period. It is the data set used by
%   Nayak et al. (2026), so results can be compared directly.
% - Period: weather years 1982-2020 (demand starts 1982-01-01). Events are identified on the continuous
%   hourly series; the 52-week Balmorel weather-year windows are not used because they omit seven
%   year-end weeks (1984, 1990, 1996, 2001, 2007, 2012, 2018) in the core of the winter season.
% - System configurations: capacity levels and interconnection bounds.
\Pending{data}

\section{Methods}
\label{sec:methods}
% Outline (01_research_plan.md section 3).
\subsection{Residual load and event definition}
\label{sec:event_definition}
% Residual load, reference supply level, cumulative deficit recursion, event features.
\Pending{event definition}

\subsection{Aggregation into weighted representative events}
\label{sec:aggregation}
\Pending{aggregation}

\subsection{Validation}
\label{sec:validation}
% Four levels; statistics matrix between the reduced and the original data set; bootstrap tolerance.
\Pending{validation}

\subsection{Representativeness under climate change}
\label{sec:climate}
% Time of emergence of the distance between the historical set and future windows; update rule.
\Pending{climate change}

\section{Results}
\label{sec:results}
% One subsection per research question (Q1-Q4); numbers only through table/text_values.tex.
\Pending{results}

\section{Discussion}
\label{sec:discussion}
\Pending{discussion}

\section{Conclusions}
\label{sec:conclusions}
\Pending{conclusions}


\section*{Acknowledgements}
% TODO: funding and data providers.
\Pending{acknowledgements}

\section*{Declaration of competing interest}
The authors declare that they have no known competing financial interests or personal
relationships that could have appeared to influence the work reported in this paper.

\section*{Data availability}
% TODO: DOIs of the input data (DTU Data doi:10.11583/DTU.c.7964141; Mopo, PECD 4.2) and of the
% event catalogue and representative event set produced by this study.
\Pending{data availability}

\bibliographystyle{elsarticle-num-names}
\bibliography{references}

\end{document}
